\section*{第 \ref{ch:b3:group-theory-applied} 章——群论应用}

\begin{solution}{exo:b3:group-theory-applied:1}
$n_{A_1} = \frac14(5 - 1 + 3 + 1) = 2$, $n_{A_2} = \frac14(5 - 1 - 3 - 1) = 0$,
$n_{B_1} = \frac14(5 + 1 + 3 - 1) = 2$, $n_{B_2} = \frac14(5 + 1 - 3 + 1) = 1$:
$\Gamma = 2A_1 + 2B_1 + B_2$，维数为 5。
\end{solution}

\begin{solution}{exo:b3:group-theory-applied:2}
与水相同，按不动原子计数得 $\Gamma_{\mathrm{vib}} = 2A_1 + B_1$。$A_1$（$z$）与
$B_1$（$x$）均具红外活性，两者亦具拉曼活性（$x^2$，$xz$）：
每张谱图中各有三条谱带。
\end{solution}

\begin{solution}{exo:b3:group-theory-applied:3}
$B_1\otimes B_2$：特征标 $(1, 1, -1, -1)$ $= A_2$。$A_2\otimes B_1$：$(1, -1, -1,
1) = B_2$。在 $C_{3v}$ 中，$E\otimes E$ 的特征标为 $(4, 1, 0)$；约化：$A_1 + A_2 +
E$.
\end{solution}

\begin{solution}{exo:b3:group-theory-applied:4}
红外：两个 $t_2$ 简正模式，3019 与 \qty{1306}{cm^{-1}}。拉曼：四个全部（$a_1$、
$e$、$t_2$ 均按二次函数变换）。无反演中心，
故允许出现重合。
\end{solution}

\begin{solution}{exo:b3:group-theory-applied:5}
不动原子数 $(4, 1, 2, 4, 1, 2)$ 乘以 $\chi_{xyz} = (3, 0, -1, 1, -2, 1)$，得
$\chi_{3N} = (12, 0, -2, 4, -2, 2)$，约化为 $A_1' + A_2' + 3E' + 2A_2'' +
E''$。去除 $E' + A_2''$（平动）与 $A_2' + E''$（转动）：$A_1' + 2E'
+ A_2''$。红外：$2E' + A_2''$，三条谱带；拉曼：$A_1' + 2E'$，三条谱带。$A_2''$
（面外弯曲）仅具红外活性，$A_1'$（对称伸缩）仅具拉曼活性，$E'$
简正模式在两种谱图中均出现。
\end{solution}

\begin{solution}{exo:b3:group-theory-applied:6}
红外：两个 $T_{1u}$ 简正模式，两条谱带。拉曼：$A_{1g}$、$E_g$、$T_{2g}$，三条谱带。
$T_{2u}$ 在两种谱中均不出现。\ce{SF6} 有反演中心，适用互斥规则。
\end{solution}

\begin{solution}{exo:b3:group-theory-applied:7}
在 $E$、$C_3$、$C_3^2$ 下，函数 $h_2$ 变为 $h_2$、$h_3$、$h_1$（特征标分别为 2，
$-1$，$-1$）；反映的贡献为 0。$\hat P_Eh_2 \propto 2h_2 - h_3 - h_1$。
其与 $2h_1 - h_2 - h_3$ 的重叠（忽略原子间重叠）为 $-3$，后者的
范数为 $\sqrt6$；减去投影后，得 $(2h_2 - h_1 - h_3) +
\frac12(2h_1 - h_2 - h_3) = \frac32(h_2 - h_3)$.
\end{solution}

\begin{solution}{exo:b3:group-theory-applied:8}
面式（$C_{3v}$，特征标 $3, 0, 1$）：$A_1 + E$，两条红外谱带。经式（$C_{2v}$，
特征标 $3, 1, 3, 1$）：$2A_1 + B_1$，三条谱带。\ce{M(CO)6}（$O_h$）：$A_{1g} + E_g
+ T_{1u}$，一条红外谱带（$T_{1u}$）。
\end{solution}

\begin{solution}{exo:b3:group-theory-applied:9}
特征标 $6, 0, 0, 2, 2, 0, 0, 0, 4, 2$：仅 $E$、沿轴的 $C_4$ 与
$C_2$（两个配体留在原位）、三个 $\sigma_h$（四个）与六个 $\sigma_d$
（两个）保留配体位置不变。约化结果为 $A_{1g} + E_g + T_{1u}$。金属 $s$
（$a_{1g}$）、$d_{z^2}$ 与 $d_{x^2-y^2}$（$e_g$）以及 $p_x$、$p_y$、$p_z$（$t_{1u}$）组合；
$t_{2g}$ 保持为非键轨道。
\end{solution}

\begin{solution}{exo:b3:group-theory-applied:10}
$n_i = \frac1h(1\cdot h\cdot\chi_i(E)) = d_i$。正则
表示的维数为 $h$，并且 $\sum_in_id_i = \sum_id_i^2$。
\end{solution}

\begin{solution}{exo:b3:group-theory-applied:11}
在 $D_{2h}$ 中（分子沿 $z$ 轴），不动原子数为 $(4, 4, 0, 0, 0, 0, 4, 4)$，$\chi_{3N}
= (12, -4, 0, 0, 0, 0, 4, 4)$，约化为 $2A_g + 2B_{2g} + 2B_{3g} + 2B_{1u} + 2B_{2u}
+ 2B_{3u}$。去除平动 $B_{1u} + B_{2u} + B_{3u}$ 及两个转动
$B_{2g} + B_{3g}$：$2A_g + B_{2g} + B_{3g} + B_{1u} + B_{2u} + B_{3u}$，共七个简正模式。
在 $D_{\infty h}$ 中：$2\Sigma_g^+ + \Pi_g + \Sigma_u^+ + \Pi_u$。红外：$\Sigma_u^+$ 与 $\Pi_u$
（两条谱带）；拉曼：$2\Sigma_g^+$ 与 $\Pi_g$（三条谱带）；无谱带重合。
\end{solution}

\begin{solution}{exo:b3:group-theory-applied:12}
$1b_2 \to 4a_1$：$B_2$，$y$，允许，垂直于平面。$3a_1 \to 4a_1$：
$A_1$，$z$，允许，沿 $C_2$ 轴。$1b_2 \to 2b_1$：$A_2$，禁阻。
$1b_1 \to 4a_1$：$B_1$，$x$，允许，位于垂直于该轴的平面内。
\end{solution}

\begin{solution}{pb:b3:group-theory-applied:1}
\textbf{1.} 顺式：两个 CO 位于相邻顶点；反式：相对；面式：三个
CO 位于同一面上；经式：三个 CO 位于同一经线上（两个互为反式，一个居于两者之间）。
\textbf{2.} $C_{2v}$.
\textbf{3.} $D_{4h}$.
\textbf{4.} $C_{3v}$.
\textbf{5.} $C_{2v}$.
\textbf{6.} 仅反式-\ce{ML4(CO)2}。
\textbf{7.} 若某操作将一个 CO 移到另一个 CO 的位置上，则其矢量偏离
对角线（贡献为 0）；保持某一 CO 不动的操作，将其键矢量
映射到自身（贡献为 +1）；无一操作使 C--O 矢量反向。
\textbf{8.} 特征标 $(2, 0, 2, 0)$，其中 $\sigma_v(xz)$ 为两个 CO 所在平面：$A_1 +
B_1$.
\textbf{9.} 特征标在 $E$、$C_4$、$C_2$、$\sigma_v$、$\sigma_d$ 下为 2，其余为 0：
$A_{1g} + A_{2u}$.
\textbf{10.} $(3, 0, 1)$: $A_1 + E$.
\textbf{11.} $(3, 1, 3, 1)$: $2A_1 + B_1$.
\textbf{12.} 2、2、3、3：每个 CO 对应一维。
\textbf{13.} 顺式 2，反式 1，面式 2，经式 3。
\textbf{14.} 顺式 2，反式 1，面式 2，经式 3。
\textbf{15.} 反式：其红外谱带（$A_{2u}$）与拉曼谱带（$A_{1g}$）分属
不同的简正模式。
\textbf{16.} $A_{2u}$ 在 $i$ 作用下改变符号：一个 CO 伸长，另一个
缩短，两者反相。
\textbf{17.} $E$ 为一对简并模式：两个频率相同的简正模式给出一条
谱带，另加 $A_1$ 谱带。
\textbf{18.} 一种反式二取代、具反演中心的异构体。
\textbf{19.} 三条红外谱带：为经式（面式将给出两条）。
\textbf{20.} 是：经式有三个拉曼活性的 CO 简正模式，频率相同，
由于不具反演中心；面式只显示两条。
\textbf{21.} $A_1$：同相伸缩是全对称的。
\textbf{22.} 面式/经式混合物至多会显示五条红外谱带，其中面式
谱带位于不同位置；三条清晰的谱带与单一异构体相符。
${}^{13}$C 核磁共振谱可加以判定：经式的两个羰基信号比为
2:1，面式为一个信号。
\textbf{23.} \textbf{经式异构体具有三个红外活性的 CO 伸缩振动，面式
异构体仅两个}：生成物为经式。
\end{solution}
